% Figure: post-buckling pitchfork bifurcation. Load P/Pcr vs midspan
% lateral deflection. Straight branch up to Pcr, then two symmetric
% deflected branches rising slightly (supercritical pitchfork).
\begin{figure}[htbp]
\centering
\begin{tikzpicture}
\begin{axis}[
  width=11cm, height=7cm,
  xlabel={midspan lateral deflection $\delta/L$},
  ylabel={axial load $P/P_{\mathrm{cr}}$},
  xmin=-0.4, xmax=0.4, ymin=0, ymax=1.6,
  axis lines=middle,
  tick label style={font=\scriptsize}, label style={font=\small},
  ytick={1.0}, yticklabels={$1$},
  xtick={0},
]
% trivial (straight) branch, stable below Pcr, unstable above
\addplot[engblue, very thick, domain=0:1.0, samples=2] coordinates {(0,0) (0,1.0)};
\addplot[engblue, very thick, dashed, domain=1.0:1.6, samples=2] coordinates {(0,1.0) (0,1.6)};
% post-buckling branches: P/Pcr = 1 + (pi^2/8)(delta/L)^2 approx
\addplot[engred, very thick, domain=0:0.38, samples=60]
  {1 + (pi^2/8)*x^2};
\addplot[engred, very thick, domain=-0.38:0, samples=60]
  {1 + (pi^2/8)*x^2};
% bifurcation point
\addplot[engblack, only marks, mark=*, mark options={fill=engamber}]
  coordinates {(0,1.0)};
\node[engamber, font=\scriptsize, anchor=south west] at (axis cs:0.02,1.02)
  {bifurcation $P=P_{\mathrm{cr}}$};
\node[engblue, font=\scriptsize, anchor=east] at (axis cs:-0.01,0.5) {straight (stable)};
\node[engblue, font=\scriptsize, anchor=west] at (axis cs:0.01,1.35) {straight (unstable)};
\node[engred, font=\scriptsize, anchor=south west] at (axis cs:0.2,1.18) {buckled branch};
\end{axis}
\end{tikzpicture}
\caption[Supercritical pitchfork at the buckling load]{Load versus
lateral deflection for a pinned-pinned column, showing the
supercritical pitchfork bifurcation at the Euler load. Below
$P_{\mathrm{cr}}$ the only equilibrium is the straight column (blue,
stable). At $P=P_{\mathrm{cr}}$ the straight state loses stability
(blue dashed) and two symmetric deflected branches appear (red),
along which the column carries slightly more than $P_{\mathrm{cr}}$
as the deflection grows. Linear Euler theory predicts the
bifurcation load but not the post-buckling path; the gentle rise of
the red branches is a result of the nonlinear large-deflection
(elastica) theory. Real columns, with imperfections, follow a
rounded curve that approaches the bifurcation diagram from one side
and never reaches the ideal straight branch.}
\label{fig:vol03ch08:postbuckling}
\end{figure}
